% from ega0-21-fr.tex
\oldpage[0\textsubscript{IV-1}]{156}
D'après (21.1.5, (iii)), $\mathfrak J_{B/A}$ est un nilidéal~; en outre,
en vertu de (20.4.4) et de l'hypothèse, $\mathfrak J_{B/A}$ est un idéal de
type fini, donc il est nilpotent~; comme la relation $\Omega^1_{B/A}=0$
signifie que $\mathfrak J_{B/A}=\mathfrak J_{B/A}^2$, on en conclut que
$\mathfrak J_{B/A}=0$, ou encore que $\pi:B\otimes_A B\to B$ est bijectif.
Par ailleurs, tout élément de $B$ ayant sa puissance $p^s$-ième dans $u(A)$,
$B$ est entier sur $A$, et étant de type fini, c'est une $A$-algèbre finie.
On est ainsi ramené à prouver le lemme suivant (où l'on ne suppose pas que
les anneaux considérés sont de caractéristique $p$) :
\end{proof}

\begin{lemma}[21.1.6.1]
\label{0.21.1.6.1-fr}
Soient $R$ un anneau, $S$ une $R$-algèbre finie~; si l'homomorphisme
canonique $\pi:S\otimes_R S\to S$ est bijectif, alors l'homomorphisme
structural $u:R\to S$ est surjectif.
\end{lemma}

\begin{proof}
Il suffit de montrer que, pour tout idéal maximal $\mathfrak m$ de $R$, si
l'on pose $T=R-\mathfrak m$, l'homomorphisme
$u_{\mathfrak m}:R_{\mathfrak m}\to T^{-1}S$ est surjectif (Bourbaki,
\emph{Alg. comm.}, chap.~II, §~3, n$^\circ$~3, th.~1)~; or l'hypothèse
entraîne que l'homomorphisme
$(T^{-1}S)\otimes_{R_{\mathfrak m}}(S/\mathfrak mS)\to T^{-1}S$
est bijectif (0$_I$, 1.3.4), et comme $T^{-1}S$ est une
$R_{\mathfrak m}$-algèbre finie, on voit qu'on peut se borner au cas où
$R$ est un anneau local.

Désignant encore par $\mathfrak m$ son idéal maximal, il suffit de prouver que
$u\otimes1:R/\mathfrak m\to S/\mathfrak mS$ est surjectif, en vertu du
lemme de Nakayama ($S$ étant un $R$-module de type fini)~; comme
l'homomorphisme canonique
$(S/\mathfrak mS)\otimes_{R/\mathfrak m}(S/\mathfrak mS)\to S/\mathfrak mS$
est bijectif et que $S/\mathfrak mS$ est une $(R/\mathfrak m)$-algèbre finie,
on est finalement ramené au cas où $R$ est un corps~; mais alors les rangs de
$S\otimes_R S$ et de $S$ sur $R$ ne peuvent être égaux que si $S$ est de rang
$0$ ou $1$, donc si $u:R\to S$ est surjectif.
\end{proof}

\begin{proposition}[21.1.7]
\label{0.21.1.7-fr}
Soient $A$ un anneau, $B$ une $A$-algèbre, $(x_\alpha)_{\alpha\in I}$ une
famille d'éléments de $B$. Considérons les propriétés suivantes~:
\begin{enumerate}
  \item[a)] $B=A[B^p,(x_\alpha)]$, autrement dit la $A[B^p]$-algèbre
  $B$ est engendrée par la famille $(x_\alpha)_{\alpha\in I}$.
  \item[b)] Le $A[B^p]$-module $B$ est engendré par les monômes
  $\prod_\alpha x_\alpha^{n(\alpha)}$, où $(n(\alpha))_{\alpha\in I}$ est
  une famille d'entiers à support fini telle que $0\leq n(\alpha)<p$ pour
  tout $\alpha\in I$.
  \item[c)] Le $B$-module $\Omega^1_{B/A}$ est engendré par les
  $d_{B/A}(x_\alpha)$ $(\alpha\in I)$.
\end{enumerate}
Alors les propriétés a) et b) sont équivalentes et entraînent c)~; si, en
outre, $B$ est une $A[B^p]$-algèbre de type fini, c) est équivalente à a) et b).
\end{proposition}

\begin{proof}
Il est clair que b) entraîne a), et inversement a) entraîne b), car tout
monôme $\prod_\alpha x_\alpha^{m(\alpha)}$, où les $m(\alpha)$ sont des
entiers $\geq0$ (formant une famille à support fini), peut s'écrire
$(\prod_\alpha x_\alpha^{q(\alpha)})^p\prod_\alpha x_\alpha^{r(\alpha)}$,
en divisant chaque $m(\alpha)$ par $p$, ce qui donne
$m(\alpha)=p q(\alpha)+r(\alpha)$ avec $0\leq r(\alpha)<p$.
Le fait que a) entraîne c) résulte de (21.1.5, (ii)) et de (20.4.7).

Supposons inversement c) vérifiée et que $B$ soit une $A[B^p]$-algèbre de
type fini~; soit $B'$ la sous-$A[B^p]$-algèbre de $B$ engendrée par les
$x_\alpha$~; dans la suite exacte (20.5.7.1)
\[
\Omega^1_{B'/A[B^p]}\otimes_{B'}B\longrightarrow
\Omega^1_{B/A[B^p]}\longrightarrow\Omega^1_{B/B'}\longrightarrow0
\]
l'hypothèse entraîne que la flèche de gauche est surjective (compte tenu de
(21.1.5, (ii)))~; on a donc $\Omega^1_{B/B'}=0$, et comme
$B^p\subset B'\subset B$ et que $B$ est une $B'$-algèbre de type fini,
on a nécessairement $B'=B$ par (21.1.6).
\end{proof}

\begin{remark}[21.1.8]
\label{0.21.1.8-fr}
Lorsque $B$ est un corps, nous démontrerons l'équivalence des propriétés a),
b) et c) sans hypothèse de finitude (21.4.5).
\end{remark}

\begin{definition}[21.1.9]
\label{0.21.1.9-fr}
Soient $A$ un anneau (de caractéristique $p$), $B$ une $A$-algèbre. On dit
qu'une famille $(x_\alpha)_{\alpha\in I}$ d'éléments de $B$ est $p$-libre sur
$A$ (resp. un système de $p$-générateurs de $B$ sur $A$, resp. une $p$-base
de $B$ sur $A$) si la famille des monômes
$\prod_\alpha x_\alpha^{n(\alpha)}$ ($0\leq n(\alpha)<p$,
$(n(\alpha))_{\alpha\in I}$ à support fini) est une famille libre (resp. un
système de générateurs, resp. une base) dans le $A[B^p]$-module $B$.
\end{definition}

On a des définitions correspondantes pour un ensemble $M$ d'éléments de $B$,
en considérant la famille définie par l'injection canonique $M\to B$.
Lorsqu'on prend pour $A$ le corps premier $\mathbf F_p$ (auquel cas
$A[B^p]=B^p$), on omet la mention de $A$ dans la définition précédente (ou
l'on dit encore qu'une famille est «absolument» $p$-libre, resp. un système
«absolu» de $p$-générateurs, resp. une $p$-base «absolue»).

